\documentclass[11pt]{article}

\usepackage[margin=1in]{geometry}
\usepackage{amsmath,amssymb,amsthm,amsfonts}
\usepackage{mathtools}
\usepackage{booktabs}
\usepackage{array}
\usepackage{tabularx}
\usepackage{xcolor}
\usepackage{hyperref}
\usepackage{algorithm}
\usepackage{algpseudocode}
\usepackage{tcolorbox}
\tcbuselibrary{skins,breakable}
\usepackage{cleveref}
\usepackage{microtype}
\usepackage{enumitem}
\usepackage{caption}
\usepackage{makecell}
\usepackage{setspace}
\onehalfspacing

% --- Theorem environments ---
\theoremstyle{plain}
\newtheorem{theorem}{Theorem}[section]
\newtheorem{proposition}[theorem]{Proposition}
\newtheorem{lemma}[theorem]{Lemma}
\newtheorem{corollary}[theorem]{Corollary}
\theoremstyle{definition}
\newtheorem{definition}[theorem]{Definition}
\newtheorem{remark}[theorem]{Remark}
\newtheorem{assumption}[theorem]{Assumption}

% --- tcolorbox styles ---
\tcolorboxenvironment{theorem}{
  enhanced,breakable,colback=blue!3!white,colframe=blue!35!black,
  boxrule=0.5pt,arc=4pt,left=6pt,right=6pt,top=3pt,bottom=3pt}
\tcolorboxenvironment{proposition}{
  enhanced,breakable,colback=blue!2!white,colframe=blue!25!black,
  boxrule=0.4pt,arc=4pt,left=6pt,right=6pt,top=3pt,bottom=3pt}
\tcolorboxenvironment{definition}{
  enhanced,colback=gray!5!white,colframe=gray!30!white,
  boxrule=0.4pt,arc=4pt,left=6pt,right=6pt}
\tcolorboxenvironment{corollary}{
  enhanced,colback=teal!3!white,colframe=teal!35!black,
  boxrule=0.4pt,arc=4pt,left=6pt,right=6pt}
\newenvironment{insightbox}[1][Insight]{%
  \begin{tcolorbox}[enhanced,breakable,colback=orange!5!white,
    colframe=orange!60!black,fonttitle=\bfseries,title=#1,
    boxrule=0.6pt,arc=5pt,left=8pt,right=8pt]}
  {\end{tcolorbox}}

% --- Math macros ---
\newcommand{\R}{\mathbb{R}}
\newcommand{\E}{\mathbb{E}}
\newcommand{\cA}{\mathcal{A}}
\newcommand{\cL}{\mathcal{L}}
\newcommand{\cN}{\mathcal{N}}
\newcommand{\cW}{\mathcal{W}}
\newcommand{\norm}[1]{\left\lVert#1\right\rVert}
\newcommand{\vnet}{v_\phi}
\newcommand{\acurr}{a_{\mathrm{curr}}}
\newcommand{\tdiff}{t}
\newcommand{\texec}{\tau}
\newcommand{\bz}{z}
\newcommand{\chunk}{\mathbf{v}}
\newcommand{\chunknoisy}{\chunk_\tdiff}
\newcommand{\chunkclean}{\chunk_0}
\newcommand{\ccoeff}{\mathbf{c}}
\newcommand{\cclean}{\ccoeff^*}
\newcommand{\cnoisy}{\ccoeff_\tdiff}
\DeclareMathOperator*{\argmin}{arg\,min}
\DeclareMathOperator{\Law}{Law}
\hypersetup{colorlinks=true,linkcolor=blue!55!black,
  citecolor=green!45!black,urlcolor=blue!55!black}

% --- Title ---
\title{
  \textbf{MOTIF: Modulated ODE Trajectory with Integrated Flow}\\[6pt]
  \large From Action Images to Action Videos:\\
  Physical-Time-Aware Robot Control via Functional Flow Matching\\
  over Fourier Coefficient Space\\[4pt]
  \normalsize\textit{Two Orthogonal Time Axes, Structural $C^\infty$ Continuity,
  Direct Velocity Supervision, and Closed-Loop Execution without Architectural Overhead}
}
\author{}
\date{\today}

\begin{document}
\maketitle
\tableofcontents
\newpage

% =============================================================================
\section{Introduction}
\label{sec:intro}
% =============================================================================

\subsection{The Action Image Problem}

Consider how $\pi_0$~\cite{black2024pi0} generates robot actions.
The action expert receives a noisy tensor
$\mathbf{a}_\tdiff \in \R^{K \times d}$---a chunk of $K$ joint
positions---and progressively denoises it via Flow Matching toward a
clean action sequence.
The integer $K = 50$ covers one second of motion at 50\,Hz;
both $K$ and the control frequency are fixed at training time.

This formulation makes a subtle but consequential error:
\textbf{it treats a sequence of robot motions as if it were a single image.}

A natural image is a function of spatial pixel coordinates $(x, y)$.
A video is a function of both spatial coordinates \emph{and physical time}
$\texec \in [0,T]$, where $\texec$ is measured in seconds.
The distinction is not cosmetic: in a video, the meaning of frame $k$
depends on \emph{when} it occurs, not merely on its ordinal rank.

In $\pi_0$, the action chunk $\mathbf{a} \in \R^{K \times d}$ is indexed by
the integer $k \in \{0,\ldots,K-1\}$.
This index encodes only \emph{order}, not physical time.
The model has no internal representation of:
\begin{enumerate}[leftmargin=2em,noitemsep]
  \item \emph{How fast} the motion unfolds---whether $a_k$ and $a_{k+1}$
    are separated by 20\,ms or 200\,ms.
  \item \emph{What velocity} the robot should have at each moment---the
    direct signal that low-level impedance controllers require.
  \item \emph{Continuous smoothness}---the discrete output provides no
    structural guarantee that neighboring actions are physically consistent.
  \item \emph{The robot's actual state}---the chunk is executed open-loop
    with no interface through which the true proprioceptive state
    can re-enter within a chunk.
\end{enumerate}

\begin{insightbox}[The Core Distinction]
\centering
\setlength{\tabcolsep}{5pt}
\begin{tabular}{lcc}
\toprule
& \textbf{$\pi_0$ (action image)} & \textbf{MOTIF (action video)} \\
\midrule
Output object       & position tensor $[K, d]$        & velocity field $v(\texec) \in C^\infty$ \\
Time axis           & ordinal step $k$ (integer)       & physical time $\texec$ (seconds) \\
Primary quantity    & position                          & velocity \\
Continuity          & none (discrete points)            & structural $C^\infty$ (Fourier basis) \\
State interface     & none (open-loop in chunk)         & $A$ is a direct query at $\tdiff=0$ \\
Control frequency   & fixed at training                 & free at inference \\
High-freq.\ data    & downsampled to $K$                & fully utilized \\
\bottomrule
\end{tabular}
\end{insightbox}

\subsection{The Right Object: A Smooth Velocity Field}

We propose that the correct output of a robot action expert is not a
position sequence but a \emph{continuous velocity field}:
\begin{equation}
  \label{eq:vfield}
  v : [0,T] \to \R^d, \qquad v(\texec \mid \bz,\, A_{\mathrm{curr}})
  = \text{velocity the robot should have at time } \texec,
\end{equation}
conditioned on the visual-linguistic context $\bz = f_{\mathrm{VLM}}(o,\ell)$
and the robot's current state $A_{\mathrm{curr}}$.

This reframing has four structural consequences:

\paragraph{Velocity as the primary output.}
$v(\texec)$ is the direct reference signal for impedance controllers.
No finite-differencing, no post-hoc smoothing.

\paragraph{Physical time as an explicit input.}
$\texec$ is measured in seconds.
The model learns that $v(0.1\,\text{s})$ and $v(0.9\,\text{s})$ are
physically different moments.
Control frequency is a free inference-time parameter.

\paragraph{Structural $C^\infty$ continuity.}
By parameterizing $v$ as a finite sum of Fourier (cosine) basis functions,
continuity is a mathematical property of the representation, not a soft
constraint imposed by the loss function.
No regularization hyperparameters are introduced.

\paragraph{State as a runtime query.}
Because the current state $A$ appears as an input to the velocity
evaluation, the field can be queried at the robot's true proprioceptive
state at every control step.
Closed-loop control is structural, not engineered.

\subsection{MOTIF: Four Changes to $\pi_0$}

We propose \textbf{MOTIF} (\textbf{M}odulated \textbf{O}DE \textbf{T}rajectory
with \textbf{I}ntegrated \textbf{F}low).
The action expert in $\pi_0$ already implements Flow Matching---an ODE
in denoising time $\tdiff$.
MOTIF does not add a new ODE; it changes \emph{what the FM denoises}:
from a position sequence (action image) to a Fourier coefficient vector
(compact representation of a $C^\infty$ velocity field).

The four changes to $\pi_0$'s action expert are:
\begin{enumerate}[leftmargin=2em,noitemsep]
  \item Replace step-index positional encoding with
    \emph{sinusoidal embedding of physical time} $\texec_k$ (in seconds).
  \item Change the FM generation space from $\R^{K \times d}$ (position
    sequence) to $\R^{(M+1) \times d}$ (DCT Fourier coefficients),
    with $M \ll K$.
  \item Add the \emph{current robot state} $A$ as a runtime query input,
    active only at execution time ($\tdiff = 0$).
  \item Supervise with \emph{velocity targets} $\dot{a}_k$ from
    high-frequency proprioceptive data, via Fourier decoding.
\end{enumerate}

The VLM backbone, training infrastructure, and data pipeline
are inherited unchanged from $\pi_0$ / openpi~\cite{openpi2025}.

\paragraph{LEITMOTIF.}
Named after the musical concept of a recurring theme that continuously
pulls a composition back to its center, \textbf{LEITMOTIF} is the
closed-loop execution mode of MOTIF.
Because the velocity field accepts the current state $A$ as a query
input, it can be evaluated at the robot's true proprioceptive state
at every control step---continuously correcting deviations.

\subsection{Contributions}

\begin{enumerate}[leftmargin=2em]
  \item \textbf{Conceptual:} The \emph{action image vs.\ action video}
    framing, which precisely identifies four structural deficiencies of
    current VLA action experts and proposes a smooth velocity field as
    the minimal correction.
  \item \textbf{Methodological:} The MOTIF architecture: four minimal
    changes to the $\pi_0$ action expert that jointly achieve
    physical time awareness, structural $C^\infty$ continuity,
    direct velocity output, and closed-loop execution capability.
  \item \textbf{Theoretical:} Formal results on FM correctness
    (Theorem~\ref{thm:funcfm}), structural smoothness
    (Theorem~\ref{thm:smooth}), sample complexity advantage
    (Theorem~\ref{thm:sample}), and closed-loop robustness
    (Proposition~\ref{prop:bibs}).
  \item \textbf{Empirical:} Five targeted experiments including the
    mid-denoising executability test (E3) and the physical-time
    perturbation test (E1), both unique to MOTIF.
\end{enumerate}

% =============================================================================
\section{The Two-Time-Axis Framework}
\label{sec:framework}
% =============================================================================

\subsection{Two Orthogonal Kinds of Time}

Every generative model for robot control operates under two distinct,
orthogonal notions of time.

\begin{definition}[Two Orthogonal Time Axes]
\label{def:two_times}
\begin{itemize}[leftmargin=2em,noitemsep]
  \item \textbf{FM denoising time} $\tdiff \in [0,1]$: parameterizes
    the generative process from noise ($\tdiff{=}1$) to a clean
    motion intention ($\tdiff{=}0$).
    This axis lives entirely within the model's computational graph
    and has no physical interpretation.
  \item \textbf{Physical execution time} $\texec \in [0,T]$: the real
    elapsed time within one action chunk, measured in seconds.
    This axis exists independently of the model.
\end{itemize}
A framework \emph{orthogonalizes} these axes if, at every denoising
step $\tdiff$, the current generative state already defines a
physically interpretable, smooth object over $\texec$.
\end{definition}

\subsection{Why $\pi_0$ Conflates the Two Axes}

In $\pi_0$, action token $k$ receives the positional encoding
$\text{sinusoidal}(k)$, implicitly setting its physical time to
$\texec_k = k / f_{\mathrm{ctrl}}$ where $f_{\mathrm{ctrl}}$ is the
\emph{training} control frequency.
This creates a rigid coupling: changing $f_{\mathrm{ctrl}}$ at inference
changes the physical meaning of every token, requiring full retraining.
Because $k$ is a dimensionless integer, the model cannot distinguish
``50 steps at 50\,Hz (1\,s)'' from ``50 steps at 200\,Hz (0.25\,s)'';
the physical scale of time is absent.

Furthermore, at intermediate denoising steps $\tdiff \in (0,1)$, the
tensor $\mathbf{a}_\tdiff$ is a convex combination of noise and clean
positions with no physical interpretation: the two axes are conflated.

\subsection{MOTIF Under This Taxonomy}

In MOTIF, each action token at physical time $\texec_k$ receives the
embedding $\text{TimeEncode}(\texec_k)$ (in seconds).
The FM generation space is the space of Fourier coefficient vectors
$\ccoeff \in \R^{(M+1)\times d}$; at every denoising step $\tdiff$,
the current $\ccoeff_\tdiff$ can be decoded into a $C^\infty$ velocity
function over $\texec \in [0,T]$.
The two axes are fully orthogonalized.

\begin{table}[ht]
\centering
\caption{Taxonomy of VLA action generation under the two-axis framework.}
\label{tab:taxonomy}
\setlength{\tabcolsep}{5pt}
\begin{tabular}{lccccccc}
\toprule
\textbf{Method} &
\makecell{Physical\\time $\texec$} &
\makecell{Velocity\\primary} &
\makecell{Structural\\continuity} &
\makecell{State\\query} &
\makecell{Var.\\freq.} &
\makecell{Prob.\\model} &
\makecell{Axes\\orthog.} \\
\midrule
ACT / Diffusion Policy  & $\times$ & $\times$ & $\times$ & $\times$ & $\times$ & $\sim$ & $\times$ \\
$\pi_0$ / GR00T         & $\times$ & $\times$ & $\times$ & $\times$ & $\times$ & $\checkmark$ & $\times$ \\
ProDMP                  & $\times$ & $\times$ & $\checkmark$ & $\times$ & $\sim$ & $\sim$ & $\times$ \\
\textbf{MOTIF (ours)}   & $\checkmark$ & $\checkmark$ & $\checkmark$ & $\checkmark$ & $\checkmark$ & $\checkmark$ & $\checkmark$ \\
\bottomrule
\end{tabular}
\end{table}

% =============================================================================
\section{MOTIF: Method}
\label{sec:method}
% =============================================================================

\subsection{Setup and Notation}

Let $o$ denote an RGB observation, $\ell$ a language instruction, and
$\bz = f_{\mathrm{VLM}}(o,\ell) \in \R^{d_z}$ the VLM context embedding.
Let $d$ be the robot's degree of freedom and $T$ the chunk duration in seconds.
At inference with control frequency $f_{\mathrm{ctrl}}$, the chunk is
evaluated at $K = \lfloor T \cdot f_{\mathrm{ctrl}} \rfloor$ physical
times $\texec_k = kT/K$, $k = 0,\ldots,K-1$.
Ground-truth velocity $\dot{a}_k \in \R^d$ at each timestep is obtained
from high-frequency proprioceptive recordings via direct encoder readout
or finite difference.

\subsection{Change 1: Physical Time Encoding}
\label{sec:time_encoding}

The most important structural change is the replacement of the ordinal
step index $k$ with physical time $\texec_k$ as the positional signal
for each action token.

\paragraph{$\pi_0$ (step index):}
\begin{equation}
  \text{PosEncode}(k) = \text{sinusoidal}(k), \quad k \in \mathbb{Z}_{\geq 0}.
\end{equation}

\paragraph{MOTIF (physical time):}
\begin{equation}
  \label{eq:time_enc}
  \text{TimeEncode}(\texec_k) = \bigl[
    \sin(\omega_1 \texec_k),\,\cos(\omega_1 \texec_k),\,\ldots,\,
    \sin(\omega_m \texec_k),\,\cos(\omega_m \texec_k)
  \bigr] \in \R^{2m},
  \quad
  \omega_i = e^{i \cdot \log(10^4)/m},
\end{equation}
where $\texec_k = kT/K$ is in \emph{seconds}.
The model associates velocities with their physical moment in time,
independently of how many tokens $K$ span the chunk.
At inference, $K$ (equivalently $f_{\mathrm{ctrl}}$) is freely varied
by adjusting $\texec_k = kT/K$ with the desired $K$, requiring no retraining.

\begin{remark}
Physical time encoding and physical time continuity are distinct properties.
The former ensures the model knows the physical meaning of each token.
The latter ensures the generated object is itself a continuous function.
MOTIF achieves both; Change~1 addresses the former and Change~2 the latter.
\end{remark}

\subsection{Change 2: Fourier Coefficient Space}
\label{sec:fourier_space}

\subsubsection{Motivation: Why Continuity Must Be Structural}

Physical robot joints require continuous velocity ($v \in C^0$);
velocity discontinuities cause instantaneous torque jumps that damage
hardware.
In $\pi_0$, continuity is not enforced; the model produces $K$ independent
real vectors with no mathematical relationship between adjacent entries.
One could attempt to impose smoothness through a regularization loss
(e.g., finite-difference penalty), but this introduces extra hyperparameters
and provides only a soft, training-dependent guarantee.

The correct approach is to parameterize the output as a member of a
function space that is continuously smooth \emph{by construction}.
For velocity functions on $[0,T]$, the appropriate space is the Sobolev
space $H^1([0,T])$, which embeds into $C^0([0,T])$ by the Sobolev
embedding theorem:
\begin{equation}
  H^1([0,T]) \hookrightarrow C^0([0,T]).
\end{equation}
The natural orthonormal basis of $H^1([0,T])$ is the cosine (DCT-II) basis,
whose elements are $\phi_k(\texec) = \cos(k\pi\texec/T)$.
The $H^1$ inner product is diagonal in this basis, with high-frequency
coefficients incurring larger norm cost---exactly what the physics demands.

\subsubsection{Fourier Basis Parameterization}

Rather than outputting $K$ velocity vectors, the action expert outputs
$M+1$ DCT-II coefficient vectors $\hat{\ccoeff} \in \R^{(M+1)\times d}$,
with $M \ll K$.
The velocity function is recovered by decoding:
\begin{equation}
  \label{eq:decode}
  \hat{v}(\texec) = \sum_{k=0}^{M} \hat{c}_k\,\phi_k(\texec),
  \qquad
  \phi_k(\texec) = \cos\!\left(\frac{k\pi\texec}{T}\right).
\end{equation}

\begin{theorem}[Structural Smoothness]
\label{thm:smooth_fourier}
For any coefficient vector $\hat{\ccoeff} \in \R^{(M+1)\times d}$
produced by the network, the decoded velocity function
$\hat{v}(\texec) \in C^\infty([0,T],\R^d)$.
The integrated position trajectory
$\hat{a}(\texec) = a_0 + \int_0^\texec \hat{v}(s)\,ds \in C^\infty$,
and acceleration $\ddot{a}$, jerk $\dddot{a}$, and all higher derivatives
are analytically available at no additional computational cost.
\end{theorem}

\begin{proof}
The result is a finite sum of $C^\infty$ functions (cosines) scaled by
constant coefficients.
Finite linear combinations of smooth functions are smooth.
The integral and all derivatives follow from the fundamental theorem of
calculus applied to each term.
\end{proof}

\begin{corollary}[Mid-Denoising Executability]
\label{cor:middenoise}
Let $\cnoisy = (1-\tdiff)\cclean + \tdiff\xi$ be the FM interpolation
at any $\tdiff \in [0,1]$.
The velocity function decoded from $\cnoisy$ is $C^\infty$ for all $\tdiff$.
Any intermediate denoising state produces a physically executable trajectory.
\end{corollary}

This is in sharp contrast to $\pi_0$, where intermediate denoising states
are convex combinations of noise and clean positions with no physical
interpretation.

\subsubsection{Selecting $M$: Energy-Based Criterion}

The number of Fourier modes $M$ is determined from the training data
via an energy threshold:
\begin{equation}
  M = \min\!\left\{m : \frac{\sum_{k=0}^{m} \|\cclean_k\|^2}
    {\sum_{k=0}^{K/2} \|\cclean_k\|^2} \geq \epsilon\right\},
  \quad \epsilon = 0.95.
\end{equation}
For joint velocity trajectories of duration $T=1\,\text{s}$ at 50\,Hz,
this yields $M \approx 8$--$16$ for smooth reaching motions and
$M \approx 20$--$32$ for tasks with rapid direction changes.
Overestimating $M$ has no negative effect beyond a marginal increase in
output head dimensionality; underestimating introduces approximation error
in high-frequency components.

\subsubsection{Frequency-Weighted Noise Prior}

Standard Gaussian noise $\xi \sim \cN(0, I_{(M+1)\times d})$ assigns
equal variance to all Fourier modes.
This is inconsistent with the physical prior: robot velocity trajectories
are low-frequency dominated, so the high-frequency modes of the data have
small magnitude.

The noise prior that is consistent with the $H^1$ structure of the target
space is the frequency-weighted Gaussian:
\begin{equation}
  \label{eq:freq_noise}
  \xi_k \sim \cN\!\left(0,\, \frac{\sigma^2}{1+\omega_k^2}\,I_d\right),
  \qquad \omega_k = \frac{k\pi}{T},
\end{equation}
which assigns smaller variance to high-frequency modes.
This is a \emph{prior} modification, not a loss modification.
The FM interpolation path from this prior to the data is geometrically
better aligned with the data manifold, providing more uniform learning
signal across the denoising trajectory.

\subsection{Change 3: State-Conditioned Velocity Query}
\label{sec:state_query}

\subsubsection{Design Principle: Denoising and Execution Are Separate Regimes}

The key insight is that the FM denoising phase ($\tdiff \in (0,1]$) and
the execution phase ($\tdiff = 0$) have fundamentally different semantics:

\begin{itemize}[leftmargin=2em,noitemsep]
  \item \textbf{Denoising phase} ($\tdiff > 0$): The model is learning
    to recover a clean motion intention from noise.
    The robot's current state is not available (we are planning before
    execution), and conditioning on it would conflate two unrelated signals.
    The state input is masked.
  \item \textbf{Execution phase} ($\tdiff = 0$): The clean motion intention
    $\chunkclean$ is fixed. The model is being queried for the velocity
    at a specific physical time and state.
    The true proprioceptive state $A_k$ is available and should be used.
\end{itemize}

\subsubsection{Implementation: $t$-Conditioned State Masking}

The state mask is conditioned on the FM denoising time $\tdiff$, not on
a training/inference flag:
\begin{equation}
  \label{eq:state_mask}
  s_k = \underbrace{\mathbf{1}[\tdiff = 0]}_{\text{is execution?}}
    \cdot f_{\mathrm{state}}(A_k)
  + \underbrace{\mathbf{1}[\tdiff > 0]}_{\text{is denoising?}}
    \cdot s_{\mathrm{mask}},
\end{equation}
where $f_{\mathrm{state}} : \R^d \to \R^{d_h}$ is a learned linear
projection, $s_{\mathrm{mask}} \in \R^{d_h}$ is a learned mask token,
and $s_k$ is added to every action token at position $k$.

\begin{remark}[Critical Correctness Condition]
The masking must depend on $\tdiff$, not on a training/inference flag.
If the mask is applied during all of training (regardless of $\tdiff$),
then at $\tdiff = 0$ during the velocity supervision loss $\cL_{\mathrm{vel}}$,
the state is also masked---meaning the model never trains on the
(state provided, $t=0$) regime that it is evaluated on at inference.
This creates a training/inference distribution mismatch that silently
nullifies the state-conditioned velocity query.
Equation~\eqref{eq:state_mask} with $\tdiff$-conditioning avoids this.
\end{remark}

\subsubsection{Closed-Loop Execution: LEITMOTIF}

In open-loop mode, the integrated state $A_k^{\mathrm{ODE}} = A_0 + \int_0^{\texec_k} \hat{v}(s)\,ds$
is used as the query.
In LEITMOTIF mode, the true proprioceptive sensor state $A_k^{\mathrm{real}}$ is read at every
control step and used instead:
\begin{equation}
  v_k = \vnet\!\bigl(\chunkclean,\; 0,\; \bz,\; A_k^{\mathrm{real}},\; \texec_k\bigr).
\end{equation}
This activates the velocity field's state-correction capability at zero
additional computational cost: it is simply a different choice of query
argument at execution time.
No special training, no additional networks, no architectural modification.

\subsection{Change 4: Velocity Supervision via Fourier Decoding}
\label{sec:vel_supervision}

\subsubsection{Loss Structure}

The full training objective is:
\begin{equation}
  \label{eq:total_loss}
  \cL = \cL_{\mathrm{FM}} + \alpha\,\cL_{\mathrm{vel}},
  \qquad \alpha = 1.0.
\end{equation}

\paragraph{FM loss (system learns the distribution of motions).}
\begin{equation}
  \label{eq:fm_loss}
  \cL_{\mathrm{FM}} = \E_{\tdiff,\,\xi,\,\cclean}
  \Bigl[\norm{\vnet\!\bigl(\cnoisy,\,\tdiff,\,\bz\bigr)
    - \bigl(\xi - \cclean\bigr)}^2\Bigr],
\end{equation}
where the linear interpolation path is:
\begin{align}
  \cnoisy &= (1-\tdiff)\,\cclean + \tdiff\,\xi, \\
  \xi &\sim \cN\!\left(0,\;\mathrm{diag}\!\left(\frac{\sigma^2}{1+\omega_k^2}\right)
    \otimes I_d\right) \quad \text{(frequency-weighted, Eq.~\ref{eq:freq_noise})}, \\
  \tdiff &\sim \mathrm{Beta}(1.5, 1) \cdot 0.999 + 0.001
    \quad \text{(biased toward $t=1$, following $\pi_0$ practice)}.
\end{align}
The FM target is the vector field $u^* = \xi - \cclean$ (constant conditional flow).
This is structurally identical to $\pi_0$'s FM loss; the only change is
that the clean target $\cclean$ is now the DCT coefficient vector of
the velocity sequence rather than a position sequence.

\paragraph{Velocity supervision loss (system learns physical time and state).}
\begin{equation}
  \label{eq:vel_loss}
  \cL_{\mathrm{vel}} = \E_{i}\!\left[
    \frac{1}{K}\sum_{j=0}^{K-1}
    \norm{\texttt{decode}\!\bigl(\hat{\ccoeff}_0^{(i)},\,\texec_j\bigr)
      - \dot{a}_j^{(i)}}^2
  \right],
\end{equation}
where:
\begin{itemize}[leftmargin=2em,noitemsep]
  \item $\hat{\ccoeff}_0^{(i)}$ is the action expert's output at $\tdiff=0$,
    queried with the clean target $\ccoeff^{*(i)}$ as input, the true
    demonstration state $A_j^{(i)}$ via Eq.~\eqref{eq:state_mask},
    and physical times $\texec_j$,
  \item $\texttt{decode}(\hat{\ccoeff}_0, \texec_j) = \sum_{k=0}^{M} (\hat{\ccoeff}_0)_k\,\phi_k(\texec_j)$
    is the Fourier decoder (Eq.~\ref{eq:decode}),
  \item $\dot{a}_j^{(i)}$ is the ground-truth joint velocity from
    high-frequency proprioception.
\end{itemize}

This formulation is clean in three ways.
First, it avoids any state integration: there is no circular dependency
between the predicted velocity and the query state.
Second, it uses demonstration states $A_j^{(i)}$ directly from the
dataset as the state query, so the training distribution for the state
input matches its inference distribution (true proprioception).
Third, the comparison is in velocity space, not coefficient space,
so the loss is interpretable in physical units (rad/s)$^2$.

\subsubsection{High-Frequency Data Utilization}

Proprioceptive data is typically recorded at 200--500\,Hz while cameras
run at 10--30\,Hz.
$\pi_0$ must downsample actions to match the fixed horizon $K$,
discarding high-frequency kinematic information.
MOTIF directly benefits: $\cL_{\mathrm{vel}}$ is evaluated at any subset
of available timesteps; denser subsets provide richer supervision without
changing the model's output dimensionality.

\subsubsection{Data Preprocessing: DCT Coefficient Extraction}

The clean coefficient target $\cclean \in \R^{(M+1)\times d}$ is
extracted offline from demonstration velocity sequences:
\begin{equation}
  c^*_k = \frac{2}{K}\sum_{j=0}^{K-1} \dot{a}_j
  \cos\!\left(\frac{k\pi(j+0.5)}{K}\right),
  \quad k = 0,\ldots,M,
\end{equation}
using \texttt{scipy.fft.dct} with \texttt{norm='ortho'}.
The result is cached as a dataset field \texttt{gt\_coeffs} and loaded
directly during training, adding zero per-step computation.

\subsection{Network Modification Summary}

The action expert's network body---all Transformer layers, VLM cross-attention,
and denoising-time embedding---is unchanged.
The only code modification is:
\begin{equation}
  \text{Output head: } \texttt{Linear}(d_{\mathrm{model}},\;K \cdot d)
  \;\to\;
  \texttt{Linear}(d_{\mathrm{model}},\;(M{+}1) \cdot d).
\end{equation}
Because $M \ll K$ (typically $M=16$, $K=50$), the output head is
\emph{smaller} than in $\pi_0$.

\subsection{Inference Algorithm}

\begin{algorithm}[H]
\caption{MOTIF Inference}
\label{alg:infer}
\begin{algorithmic}[1]
\Require Observation $o$, language $\ell$, current state $\acurr$,
  control frequency $f_{\mathrm{ctrl}}$, FM steps $S$, chunk duration $T$,
  Fourier modes $M$,
  \texttt{mode} $\in \{\texttt{open\_loop},\;\texttt{leitmotif}\}$
\State $\bz \leftarrow f_{\mathrm{VLM}}(o, \ell)$
  \Comment{VLM forward pass; KV-cache reusable}
\State $K \leftarrow \lfloor T \cdot f_{\mathrm{ctrl}} \rfloor$
  \Comment{Number of control steps at target frequency}
\State $\ccoeff_1 \sim \cN\!\left(0,\;\mathrm{diag}(\sigma^2/(1{+}\omega_k^2)) \otimes I_d\right)$
  \Comment{Frequency-weighted noise in coefficient space}
\For{$s = S, S-1, \ldots, 1$}
  \Comment{FM denoising in coefficient space}
  \State $\tdiff \leftarrow s/S$
  \State $\ccoeff_{\tdiff - 1/S} \leftarrow \ccoeff_\tdiff
    - \tfrac{1}{S}\,\vnet(\ccoeff_\tdiff,\,\tdiff,\,\bz)$
\EndFor
\State $\cclean \leftarrow \ccoeff_0$
  \Comment{Clean coefficient vector (motion intention)}
\State $A \leftarrow \acurr$
\For{$k = 0, \ldots, K-1$}
  \State $\texec_k \leftarrow k \cdot T / K$
    \Comment{Physical time in seconds}
  \If{\texttt{mode = open\_loop}}
    \State $A_{\mathrm{query}} \leftarrow A$
  \ElsIf{\texttt{mode = leitmotif}}
    \State $A_{\mathrm{query}} \leftarrow \mathrm{proprioception}()$
      \Comment{True sensor state; LEITMOTIF closed-loop}
  \EndIf
  \State $v_k \leftarrow \vnet(\cclean,\;0,\;\bz,\;A_{\mathrm{query}},\;\texec_k)$
    \Comment{State-conditioned velocity query}
  \State \textbf{send} $v_k$ to impedance controller
  \State $A \leftarrow A + v_k / f_{\mathrm{ctrl}}$
    \Comment{ODE integration (open-loop only)}
\EndFor
\end{algorithmic}
\end{algorithm}

\paragraph{Key inference properties.}
\begin{itemize}[leftmargin=2em,noitemsep]
  \item \textbf{Variable frequency.} $f_{\mathrm{ctrl}}$ is set freely at
    inference by adjusting $K$. No retraining.
  \item \textbf{LEITMOTIF.} Closed-loop execution is a one-line mode switch.
  \item \textbf{Chunk boundary consistency.} Initial condition $A = \acurr$
    is the true sensor state; $C^0$ continuity at boundaries is guaranteed
    by construction.
  \item \textbf{Mid-denoising executability.} At any $\tdiff$, $\ccoeff_\tdiff$
    decodes to a $C^\infty$ velocity function.
    The FM denoising can be interrupted and the partial result executed
    on the robot.
\end{itemize}

% =============================================================================
\section{Theoretical Analysis}
\label{sec:theory}
% =============================================================================

\subsection{Assumptions}

\begin{assumption}[Smooth Action Expert]
\label{asm:smooth}
$\vnet(\cclean, 0, \bz, \cdot, \cdot)$ is $C^\infty$ in $(A, \texec)$
for any fixed $(\cclean, \bz)$, and $L$-Lipschitz in $A$ uniformly
over $\texec \in [0,T]$.
This holds for smooth activations (SiLU, GELU) on bounded inputs.
\end{assumption}

\begin{assumption}[Comparable Hypothesis Classes]
\label{asm:hypothesis}
The $\pi_0$ discrete position class $\mathcal{F}_K$ and MOTIF velocity
field class $\mathcal{F}_s$ share the same VLM backbone and are both
$\Lambda$-Lipschitz in the VLM embedding $\bz$.
The Sobolev index $s$ of $\mathcal{F}_s$ satisfies $s \geq 2$
(guaranteed for smooth activations; for ReLU, $s=1$).
\end{assumption}

\subsection{Theorem 1: FM Correctness over Coefficient Space}

Standard FM~\cite{lipman2022flow} learns a distribution over vectors in $\R^n$.
MOTIF's FM operates over coefficient vectors $\ccoeff \in \R^{(M+1)\times d}$,
which represent velocity fields via the Fourier decoder.

\begin{definition}[Functional Flow Matching]
\label{def:funcfm}
Let $p_{\mathrm{demo}}$ be the distribution over demonstration velocity
fields.
MOTIF's FM approximates the pushforward of $p_{\mathrm{demo}}$ via:
\begin{equation}
  \frac{d\ccoeff_\tdiff}{d\tdiff} = \vnet(\ccoeff_\tdiff, \tdiff, \bz),
  \quad
  \ccoeff_1 \sim \cN(0, \Sigma_\omega),
  \quad
  \texttt{decode}(\ccoeff_0,\,\cdot) \xrightarrow{M \to \infty} v^*(\cdot),
\end{equation}
where $\Sigma_\omega = \mathrm{diag}(\sigma^2/(1+\omega_k^2)) \otimes I_d$.
\end{definition}

\begin{theorem}[Functional FM Correctness]
\label{thm:funcfm}
Under standard regularity conditions, if $\cL_{\mathrm{FM}} \to 0$ and
$M$ is sufficiently large, then:
\begin{equation}
  \mathcal{W}_2\!\bigl(\Law(\texttt{decode}(\ccoeff_0)),\;p_{\mathrm{demo}}\bigr)
  \leq \varepsilon_{\mathrm{FM}} + \varepsilon_{\mathrm{approx}},
\end{equation}
where $\varepsilon_{\mathrm{FM}}$ is the FM approximation error and
$\varepsilon_{\mathrm{approx}} = O(M^{-s})$ is the Fourier truncation
error for velocity fields in $W^{s,\infty}$.
\end{theorem}

\begin{proof}[Proof sketch]
$\varepsilon_{\mathrm{FM}}$ follows from standard FM convergence
results~\cite{lipman2022flow}.
$\varepsilon_{\mathrm{approx}}$ follows from the rate of convergence of
finite Fourier series for $W^{s,\infty}$ functions: the approximation
error of the $M$-term truncation in $L^\infty$ is $O(M^{-s})$.
The triangle inequality in $\mathcal{W}_2$ yields the stated bound.
\end{proof}

\subsection{Theorem 2: Smoothness Along the Denoising Path}

\begin{theorem}[Denoising Path Smoothness]
\label{thm:smooth}
Under Assumption~\ref{asm:smooth}, for any $\tdiff \in [0,1]$:
\begin{enumerate}[leftmargin=2em,noitemsep]
  \item The velocity function
    $\hat{v}_\tdiff(\texec) = \texttt{decode}(\cnoisy, \texec) \in C^\infty([0,T])$.
  \item The trajectory from $\acurr$ under $\hat{v}_\tdiff$ is bounded:
    $\norm{\hat{a}_\tdiff(\texec)}_\infty \leq (\norm{\acurr} + C_v T)e^{LT}$.
  \item $\hat{v}_\tdiff$ is continuous in $\tdiff$ (the path
    $t \mapsto \cnoisy$ is linear by construction).
\end{enumerate}
\end{theorem}

\begin{proof}
For (1): $\cnoisy = (1-\tdiff)\cclean + \tdiff\xi \in \R^{(M+1)\times d}$
is a fixed coefficient vector; its Fourier decoding is a finite sum of
$C^\infty$ basis functions.
For (2): Gr\"{o}nwall applied to the ODE $\dot{\hat{a}} = \hat{v}_\tdiff(\hat{a}, \texec)$
with $L$-Lipschitz dynamics.
For (3): linearity in $\tdiff$ of $\cnoisy$ implies pointwise continuity of
the decoded velocity function.
\end{proof}

\begin{remark}[Mid-Denoising Executability as a Testable Prediction]
Theorem~\ref{thm:smooth} implies that at $\tdiff = 0.5$ (halfway through
denoising), the robot should exhibit coherent task-relevant motion, not
random noise.
This is a testable, quantitative prediction unique to MOTIF:
Experiment~E3 evaluates it directly.
\end{remark}

\subsection{Theorem 3: Sample Complexity Advantage}

\begin{theorem}[Sample Complexity Gap]
\label{thm:sample}
Under Assumptions~\ref{asm:smooth} and~\ref{asm:hypothesis}, let
$\hat{\pi}_N^{K}$ and $\hat{\pi}_N^{\mathrm{MOTIF}}$ be empirical risk
minimizers over the discrete position class $\mathcal{F}_K$ and the MOTIF
velocity field class $\mathcal{F}_s$:
\begin{align}
  \E[\mathcal{L}(\hat{\pi}_N^K)] - \mathcal{L}(\pi^*_K)
    &\leq O\!\left(N^{-\frac{2}{2+Kd}}\right), \label{eq:pi0_rate}\\
  \E[\mathcal{L}(\hat{\pi}_N^{\mathrm{MOTIF}})] - \mathcal{L}(\pi^*_{\mathrm{MOTIF}})
    &\leq O\!\left(N^{-\frac{2s}{2s+d}}\right). \label{eq:motif_rate}
\end{align}
For $K \geq 2$, $s \geq 2$, Eq.~\eqref{eq:motif_rate} is strictly tighter,
and the advantage grows monotonically with $K$.
\end{theorem}

\begin{proof}[Proof sketch]
Covering numbers: $\log \mathcal{N}(\mathcal{F}_K, \varepsilon) = O(Kd \log(1/\varepsilon))$
from standard uniform convergence;
$\log \mathcal{N}(\mathcal{F}_s, \varepsilon) = O((d/s)\log(1/\varepsilon))$
by Kolmogorov entropy of $W^{s,\infty}([0,T]\times\R^d)$.
Since $Kd > d/s$ for $Ks > 1$, the MOTIF bound is tighter.
For smooth activations, $s = \infty$ and MOTIF approaches the parametric
rate $O(N^{-1})$.
\end{proof}

The intuition: $\pi_0$ must learn an independent real vector at each of
the $K$ timesteps; MOTIF needs only to learn the $M+1 \ll K$ Fourier
modes of a smooth function.
The effective dimension of what must be learned is smaller.

\subsection{Proposition 4: Closed-Loop Robustness (LEITMOTIF)}

\begin{proposition}[Bounded Input, Bounded State Robustness]
\label{prop:bibs}
Suppose the learned velocity field is contracting in state:
\begin{equation}
  \tfrac{1}{2}\bigl(\nabla_A \vnet + (\nabla_A \vnet)^\top\bigr)
  \preceq -\lambda I \quad \forall\,A, \texec.
\end{equation}
Under LEITMOTIF execution with bounded state perturbation
$\norm{A_{\mathrm{real},k} - \hat{a}(\texec_k)} \leq \bar{d}$:
\begin{equation}
  \limsup_{k \to \infty}\norm{A_{\mathrm{real},k} - \hat{a}(\texec_k)}
  \leq \frac{\bar{d}}{\lambda}.
\end{equation}
\end{proposition}

Contraction rate $\lambda$ can be encouraged via a regularization term
$\gamma \cdot \E[\lambda_{\max}(\nabla_A \vnet + (\nabla_A \vnet)^\top)]$
added to the training loss (disabled by default; $\gamma = 0.01$ for
robustness-critical deployments).
Even without this term, LEITMOTIF introduces no training/inference
distribution shift at the state interface: demonstrations provide real
proprioceptive states as training examples, and LEITMOTIF queries real
states at inference.

% =============================================================================
\section{Related Work}
\label{sec:related}
% =============================================================================

\paragraph{Flow Matching for VLA.}
$\pi_0$~\cite{black2024pi0} applies conditional FM to robot action generation,
treating the action chunk as a continuous sample in $\R^{K \times d}$.
MOTIF replaces this with FM over Fourier coefficient vectors, changing the
target from ordinal-indexed positions to physically-timed velocities with
structural $C^\infty$ continuity.
The VLM backbone and FM infrastructure are unchanged; MOTIF is a drop-in
replacement for the action expert head.

\paragraph{FM-based video generation.}
FLOAT~\cite{ki2025float} uses FM in motion latent space with per-frame
temporal conditioning, demonstrating that separating frame-level
conditioning from the FM denoising axis improves temporal consistency.
MOTIF's per-token physical time encoding is a direct import of this
insight: each action token receives its own physical-time embedding,
independent of the FM denoising step.

\paragraph{Movement primitives.}
ProDMP~\cite{li2023prodmp} learns distributions over Dynamic Movement
Primitives, using fixed radial basis functions to represent position
trajectories.
MOTIF shares the use of a compact basis representation but operates in
velocity space, uses Fourier (cosine) rather than radial basis functions,
and learns the full distribution via Flow Matching rather than fitting
parametric distributions.
ProDMP does not encode physical time explicitly and generates positions
rather than velocities.

\paragraph{Continuous-time robot control.}
Neural ODEs~\cite{chen2018neural} parameterize dynamics as ODEs in physical
time but do not model distributions over dynamical laws.
MOTIF models a distribution over velocity fields via FM while also
organizing the learned field along physical time---achieving both
probabilistic modeling and continuous-time representation.

\paragraph{Action coherence in flow-based policies.}
ACG~\cite{acg2025} addresses temporal incoherence in FM-based VLAs via
inference-time guidance, targeting the symptom rather than the structural
cause.
MOTIF addresses the cause: ordinal-indexed positions have no physical
time structure, so incoherence is unavoidable under that parameterization.

% =============================================================================
\section{Experimental Design}
\label{sec:exp}
% =============================================================================

\subsection{Baselines}

\begin{enumerate}[leftmargin=2em,label=(\roman*)]
  \item $\pi_0$~\cite{black2024pi0} at fixed 50\,Hz (same VLM backbone).
  \item $\pi_0$ with temporal ensembling (inference-time smoothing).
  \item ProDMP~\cite{li2023prodmp} with flow matching over DMP parameters.
  \item \textbf{MOTIF (open-loop)}: all four changes; LEITMOTIF disabled.
  \item \textbf{MOTIF (LEITMOTIF)}: full closed-loop mode.
\end{enumerate}

\subsection{Benchmarks}

\begin{itemize}[leftmargin=2em,noitemsep]
  \item \textbf{LIBERO}~\cite{liu2023libero}: 10-task manipulation suite;
    data efficiency and within-distribution generalization.
  \item \textbf{FurnitureBench}~\cite{heo2023furniturebench}: long-horizon
    assembly; temporal coherence over multi-minute trajectories.
  \item \textbf{Real robot} (UR5e / Franka): peg insertion and cable routing;
    contact compliance and high-frequency velocity control.
\end{itemize}

\subsection{Targeted Experiments}

\paragraph{E1: Physical time robustness.}
Train MOTIF and $\pi_0$ at 50\,Hz.
At inference, evaluate at $\{10, 50, 200\}$\,Hz without retraining.
\textit{Hypothesis}: MOTIF maintains performance across frequencies
(physical time is in seconds, invariant to $K$);
$\pi_0$ degrades outside its training frequency.

\paragraph{E2: High-frequency data benefit.}
Train with proprioceptive data at $\{10, 50, 200, 500\}$\,Hz.
$\pi_0$ downsamples all to 50\,Hz; MOTIF uses all frequencies directly.
\textit{Hypothesis}: MOTIF improves monotonically with data frequency;
$\pi_0$ is insensitive beyond 50\,Hz.

\paragraph{E3: Mid-denoising executability.}
Stop FM denoising at $\tdiff \in \{0.5, 0.25, 0.1, 0.0\}$.
Execute the partial result directly.
\textit{Hypothesis}: MOTIF achieves $> 50\%$ task success at $\tdiff = 0.5$
(Corollary~\ref{cor:middenoise});
$\pi_0$ achieves $\approx 0\%$ (intermediate state is a noise-position mixture).

\paragraph{E4: LEITMOTIF closed-loop robustness.}
Apply controlled mid-chunk disturbances (joint torque impulses).
Compare MOTIF (LEITMOTIF) vs.\ MOTIF (open-loop) vs.\ $\pi_0$.
\textit{Hypothesis}: LEITMOTIF recovers within $\lambda^{-1}$ seconds;
$\pi_0$ does not recover within a chunk.

\paragraph{E5: Trajectory smoothness.}
Measure jerk $\norm{\dddot{a}}_2$ along executed trajectories.
\textit{Hypothesis}: MOTIF achieves significantly lower jerk than $\pi_0$
(structural $C^\infty$ vs.\ no continuity guarantee);
the gap grows with $K$.

\subsection{Ablations}

\begin{center}
\begin{tabular}{llll}
\toprule
ID & Variable & Values & Tests \\
\midrule
A1 & Time encoding & Step index $k$ vs.\ $\texec_k$ (s) & Change~1 value \\
A2 & FM space & Position $\R^{K\times d}$ vs.\ coefficients $\R^{(M+1)\times d}$ & Change~2 value \\
A3 & Noise prior & Standard Gaussian vs.\ frequency-weighted & Prior alignment \\
A4 & State query & With / without at $\tdiff=0$ & Change~3 value \\
A5 & Supervision & Position MSE vs.\ velocity via decode & Change~4 value \\
A6 & $M$ (modes) & $\{4, 8, 16, 32\}$ & Mode count sensitivity \\
A7 & FM steps $S$ & $\{1, 2, 4, 10\}$ & Few-step sufficiency \\
A8 & Contraction $\gamma$ & $\{0, 0.01, 0.1\}$ & LEITMOTIF robustness tradeoff \\
\bottomrule
\end{tabular}
\end{center}

% =============================================================================
\section{Conclusion}
\label{sec:conclusion}
% =============================================================================

We have presented MOTIF (Modulated ODE Trajectory with Integrated Flow),
which addresses a fundamental representational limitation of current VLA
action experts: they generate action images (position tensors indexed by
ordinal step) rather than action videos (smooth velocity fields indexed
by physical time).

Four minimal changes to the $\pi_0$ action expert---physical time
encoding, Fourier coefficient parameterization, state-conditioned
velocity query, and velocity supervision via Fourier decoding---jointly
achieve: physical time awareness independent of control frequency,
structural $C^\infty$ continuity guaranteed by the Fourier basis (not
by regularization), direct velocity output for impedance controllers,
and closed-loop execution (LEITMOTIF) as a structural property.

These capabilities emerge not from new architectures or training stages,
but from a change in what the model is asked to represent: not
``the $k$-th position in a sequence,'' but ``the velocity at physical
time $\texec$.''
The network body is unchanged; only the output head dimension changes.

\paragraph{Future directions.}
The Fourier parameterization enables a natural extension to adaptive
query density: generate a sparse coefficient vector, then decode at
higher density in physically critical intervals (e.g., contact transitions)
and lower density in smooth approach phases.
This would make the effective control frequency task-adaptive, concentrating
computation where it matters most.

% --- References ---
\begin{thebibliography}{99}

\bibitem{black2024pi0}
Black, K.\ et al.\ (2024).
$\pi_0$: A Vision-Language-Action Flow Model for General Robot Control.
\textit{arXiv:2410.24164}.

\bibitem{openpi2025}
Physical Intelligence.\ (2025).
openpi: Open-source models and packages for robotics.
\texttt{github.com/Physical-Intelligence/openpi}.

\bibitem{chen2018neural}
Chen, R.T.Q., Rubanova, Y., Bettencourt, J., \& Duvenaud, D.\ (2018).
Neural Ordinary Differential Equations.
\textit{NeurIPS} 31.

\bibitem{acg2025}
ACG Authors.\ (2025).
ACG: Action Coherence Guidance for Flow-based VLA Models.
\textit{arXiv:2510.22201}.

\bibitem{heo2023furniturebench}
Heo, M.\ et al.\ (2023).
FurnitureBench: Reproducible Real-World Benchmark for Long-Horizon Complex Manipulation.
\textit{RSS 2023}.

\bibitem{ki2025float}
Ki, T., Min, D., \& Chae, G.\ (2025).
FLOAT: Generative Motion Latent Flow Matching for Audio-driven Talking Portrait.
\textit{ICCV 2025}.

\bibitem{li2023prodmp}
Li, G.\ et al.\ (2023).
ProDMP: A Unified Perspective on Dynamic and Probabilistic Movement Primitives.
\textit{RA-L} 8(2).

\bibitem{lipman2022flow}
Lipman, Y., Chen, R.T.Q., Ben-Hamu, H., Nickel, M., \& Le, M.\ (2022).
Flow Matching for Generative Modeling.
\textit{ICLR 2023}.

\bibitem{liu2023libero}
Liu, B.\ et al.\ (2023).
LIBERO: Benchmarking Knowledge Transfer in Lifelong Robot Learning.
\textit{NeurIPS 2023}.

\bibitem{wainwright2019high}
Wainwright, M.J.\ (2019).
\textit{High-Dimensional Statistics}. Cambridge University Press.

\end{thebibliography}

% =============================================================================
\appendix
\renewcommand{\thesection}{\Alph{section}}
\setcounter{section}{0}
% =============================================================================

\section{Hyperparameter Reference}
\label{app:hyperparams}

\begin{center}
\begin{tabular}{llll}
\toprule
Parameter & Default & Determined by & Notes \\
\midrule
$M$ (Fourier modes) & 16 & Energy threshold on training set & Most important new parameter \\
$\alpha$ ($\cL_{\mathrm{vel}}$ weight) & 1.0 & Ablation A5 & \\
$T$ (chunk duration) & 1.0\,s & Task & In seconds \\
$S$ (FM steps) & 4 (fast) / 10 (quality) & Ablation A7 & \\
$\sigma$ (noise scale) & 1.0 & Rarely needs tuning & \\
$f_{\mathrm{ctrl}}$ (control freq.) & 50\,Hz & Inference-time choice & No retraining needed \\
$\gamma$ (contraction) & 0.0 & 0.01 for LEITMOTIF robustness & \\
Time embed dim $m$ & 64 & Architecture & Sinusoidal, physical seconds \\
\bottomrule
\end{tabular}
\end{center}

\section{Complete Implementation}
\label{app:impl}

\subsection{Data Preprocessing}

\begin{verbatim}
# preprocess/extract_coeffs.py
import scipy.fft
import numpy as np


def extract_dct_coeffs(v_star: np.ndarray, M: int) -> np.ndarray:
    """
    Extract DCT-II coefficients from a demonstration velocity sequence.
    Run offline; cache result as dataset field 'gt_coeffs'.

    Args:
        v_star: (K, d) demonstration velocity sequence
        M:      number of Fourier modes to retain

    Returns:
        (M+1, d) DCT-II coefficient array
    """
    coeffs_full = scipy.fft.dct(v_star, axis=0, norm='ortho')  # (K, d)
    return coeffs_full[:M + 1, :]


def decode_coeffs_to_velocity(
    coeffs: np.ndarray,   # (M+1, d)
    tau:    np.ndarray,   # (N_query,)  physical time in seconds
    T:      float,
    K_norm: int,          # original sequence length for ortho normalization
) -> np.ndarray:          # (N_query, d)
    """
    Decode DCT coefficients to velocity at arbitrary physical times.
    Implements Eq. (3.2) of the paper.
    """
    M = coeffs.shape[0] - 1
    k = np.arange(0, M + 1, dtype=np.float32)
    norm = np.ones(M + 1)
    norm[0] = 1.0 / np.sqrt(2)
    norm = norm * np.sqrt(2.0 / K_norm)
    basis = np.cos(np.pi * k[:, None] * tau[None, :] / T)  # (M+1, N_query)
    return (coeffs * norm[:, None]).T @ basis               # (N_query, d)


def estimate_M(dataset, energy_threshold: float = 0.95) -> int:
    """
    Compute the minimum M covering energy_threshold of mean spectral energy.
    Run once before preprocessing; fix M afterward.
    """
    cumulative_energies = []
    for v_star in dataset:
        coeffs = scipy.fft.dct(v_star, axis=0, norm='ortho')
        energy = (coeffs ** 2).sum(axis=-1)
        cumulative_energies.append(np.cumsum(energy) / energy.sum())
    mean_cumsum = np.stack(cumulative_energies).mean(axis=0)
    return int(np.argmax(mean_cumsum >= energy_threshold)) + 1
\end{verbatim}

\subsection{Noise Sampling}

\begin{verbatim}
# training/noise.py
import jax
import jax.numpy as jnp


def sample_freq_weighted_noise(
    M: int, d: int, T: float, B: int,
    sigma: float = 1.0, key=None,
) -> jnp.ndarray:
    """
    Sample frequency-weighted Gaussian noise (Eq. 3.4 of the paper).
    Replaces jax.random.normal(...) for the FM noise draw.

    Returns: (B, M+1, d)
    """
    k_idx   = jnp.arange(0, M + 1, dtype=jnp.float32)
    omega_k = k_idx * jnp.pi / T
    std_k   = sigma / jnp.sqrt(1.0 + omega_k ** 2)   # (M+1,)
    xi      = jax.random.normal(key, (B, M + 1, d))
    return xi * std_k[None, :, None]
\end{verbatim}

\subsection{Action Expert Modifications}

\begin{verbatim}
# models/action_expert.py  (changes from pi0 highlighted)

class MotifActionExpert(nn.Module):
    def __init__(self, d_model, d_action, M, time_embed_dim=64):
        super().__init__()
        # [CHANGE 1] Physical time embedding (replaces step-index lookup)
        self.time_embed  = MotifTimeEmbedding(time_embed_dim, d_model)

        # [CHANGE 2] Output head: (M+1)*d instead of K*d
        self.output_head = nn.Linear(d_model, (M + 1) * d_action)

        # [CHANGE 3] State query projection and mask token
        self.state_proj  = nn.Linear(d_action, d_model)
        self.state_mask  = nn.Parameter(torch.zeros(d_model))

        # ... (Transformer body unchanged from pi0) ...

    def forward(self, c_noisy, tau, t_diffusion, z_vlm,
                A_query=None):
        """
        c_noisy:     (B, M+1, d)  noisy coefficient vector
        tau:         (B, M+1)     physical times for M+1 tokens
        t_diffusion: (B,)         FM denoising time
        z_vlm:       (B, d_z)     VLM embedding
        A_query:     (B, d)       current robot state (optional)
        """
        B, Mp1, d = c_noisy.shape

        # [CHANGE 1] Physical time embedding
        time_emb = self.time_embed(tau)                     # (B, M+1, d_model)
        x = self.token_proj(c_noisy) + time_emb            # (B, M+1, d_model)

        # [CHANGE 3] State injection, conditioned on t_diffusion
        if A_query is not None:
            is_exec = (t_diffusion == 0).float()[:, None]  # (B, 1)
            state_emb = self.state_proj(A_query)            # (B, d_model)
            state_tok = (
                is_exec * state_emb
                + (1 - is_exec) * self.state_mask[None]
            )                                               # (B, d_model)
            x = x + state_tok.unsqueeze(1)                 # broadcast to (B, M+1, d_model)

        # ... Transformer forward pass (cross-attend to z_vlm, etc.) ...

        # [CHANGE 2] Output head: coefficients, not positions
        out = self.output_head(x)                          # (B, M+1, d*... )
        return out.reshape(B, Mp1, d)                      # (B, M+1, d)
\end{verbatim}

\subsection{Training Loss}

\begin{verbatim}
# training/motif_loss.py
import jax, jax.numpy as jnp
from typing import Dict, Tuple


def motif_loss(
    transformer, batch: Dict, z_vlm: jnp.ndarray,
    M: int, T: float = 1.0, alpha: float = 1.0,
    sigma: float = 1.0, key=None,
) -> Tuple[jnp.ndarray, Dict]:
    """
    Full MOTIF training loss: L = L_FM + alpha * L_vel.

    batch fields:
        'gt_coeffs': (B, M+1, d)  pre-computed DCT coefficients
        'gt_vel':    (B, K, d)    demonstration velocity sequence
        'tau':       (B, K)       physical times (seconds)
        'gt_states': (B, K, d)    demonstration proprioceptive states
    """
    B, K, d = batch['gt_vel'].shape
    c_star   = batch['gt_coeffs']    # (B, M+1, d)
    v_star   = batch['gt_vel']       # (B, K, d)
    tau      = batch['tau']          # (B, K)
    A_demo   = batch['gt_states']    # (B, K, d)

    key, k1, k2 = jax.random.split(key, 3)

    # -- L_FM: standard Flow Matching in coefficient space ----------
    c0  = sample_freq_weighted_noise(M, d, T, B, sigma, k1)  # noise
    t   = jax.random.beta(k2, 1.5, 1.0, (B,)) * 0.999 + 0.001
    c_t = (1 - t[:, None, None]) * c_star + t[:, None, None] * c0
    u_star = c0 - c_star             # FM vector field target

    tau_coeff = make_coeff_times(M, T, B)  # (B, M+1) times for coefficient tokens

    u_pred = transformer(c_t, tau=tau_coeff, t_diffusion=t, z_vlm=z_vlm)
    L_fm = jnp.mean((u_pred - u_star) ** 2)

    # -- L_vel: Fourier-decoded velocity supervision ---------------
    t_zero   = jnp.zeros((B,))
    # Use demonstration states (A_demo) as state query at t=0
    # Model sees (c_star, t=0, A_demo[:, k, :]) and should output
    # the decoded velocity matching v_star[:, k, :].
    # We pass A_demo in a flattened way: broadcast over K, take mean.
    # Simpler: query with the mean state, or with each step separately.
    # Here we use a per-step loop (vmapped for efficiency):
    def query_vel_step(A_k, tau_k):
        # A_k: (B, d),  tau_k: (B,)
        c_hat = transformer(
            c_star,
            tau=tau_coeff,
            t_diffusion=t_zero,
            z_vlm=z_vlm,
            A_query=A_k,
        )                            # (B, M+1, d)
        # Decode to velocity at tau_k
        # phi_k(tau) = cos(k*pi*tau/T)
        k_idx = jnp.arange(0, M + 1, dtype=jnp.float32)
        basis = jnp.cos(jnp.pi * k_idx * tau_k[:, None] / T)  # (B, M+1)
        return jnp.einsum('bmd,bm->bd', c_hat, basis)          # (B, d)

    # vmap over K timesteps
    A_steps   = jnp.transpose(A_demo, (1, 0, 2))   # (K, B, d)
    tau_steps = jnp.transpose(tau,    (1, 0))        # (K, B)
    v_decoded = jax.vmap(query_vel_step)(A_steps, tau_steps)  # (K, B, d)
    v_decoded = jnp.transpose(v_decoded, (1, 0, 2))           # (B, K, d)

    L_vel = jnp.mean((v_decoded - v_star) ** 2)

    total = L_fm + alpha * L_vel
    return total, {'L_FM': float(L_fm), 'L_vel': float(L_vel), 'total': float(total)}
\end{verbatim}

\subsection{Validation Checklist}

\begin{enumerate}[leftmargin=2em,label=\textbf{V\arabic*.}]
  \item $\cL_{\mathrm{vel}} < 0.01\,(\text{rad/s})^2$ within 50k steps.
    If stalled above 0.05, verify that \texttt{gt\_vel} is computed from
    high-frequency proprioception, not finite-differenced positions.

  \item $\cL_{\mathrm{FM}}$ decreases monotonically; final value comparable
    to $\pi_0$'s FM loss.
    If diverging, check that \texttt{gt\_coeffs} is zero-mean unit-variance
    normalized per dimension.

  \item \textbf{Frequency robustness pre-screen (E1).}
    After 50k steps, evaluate at 10\,Hz, 50\,Hz, 200\,Hz.
    Trajectory shapes should be qualitatively similar.
    If 200\,Hz is erratic, verify that $\texec_k$ is in seconds, not steps.

  \item \textbf{Mid-denoising sanity (E3 pre-screen).}
    Execute $\ccoeff_{t=0.5}$ (stop denoising halfway).
    The robot should approach the object coherently.
    If stationary or diverging, the velocity supervision loss is not
    propagating gradients to the time embedding.

  \item \textbf{State mask correctness.}
    At training time with $\tdiff = 0$, confirm the action expert receives
    \texttt{A\_demo} (real state), not \texttt{state\_mask}.
    At $\tdiff > 0$, confirm it receives \texttt{state\_mask}.
    Violation is the most common silent bug: the state query is active
    during denoising and inactive during velocity supervision,
    which nullifies Change~3.

  \item \textbf{LEITMOTIF chunk boundary.}
    Over 10 consecutive chunks, measure
    $\norm{a_0^{(k+1)} - a_{K-1}^{(k)}}_\infty$.
    Should be machine-precision zero (initial condition is always true sensor state).
\end{enumerate}

\subsection{Common Pitfalls}

\begin{description}[leftmargin=2em]
  \item[\textbf{P1: Time in wrong units.}]
    Passing step indices $k$ instead of physical times $\texec_k = kT/K$ to
    \texttt{MotifTimeEmbedding}.
    \textit{Symptom}: E1 fails; 200\,Hz and 50\,Hz trajectories look completely
    different.
    \textit{Fix}: always compute
    \texttt{tau = torch.arange(K) * T / K}.

  \item[\textbf{P2: Velocity not normalized.}]
    The FM target \texttt{gt\_coeffs} must be normalized identically to how
    $\pi_0$ normalizes position chunks.
    \textit{Symptom}: $\cL_{\mathrm{FM}}$ is $10\times$ the $\pi_0$ baseline.
    \textit{Fix}: normalize per-dimension mean and std of
    $\{\dot{a}_k\}$ across the training set.

  \item[\textbf{P3: State mask conditioned on \texttt{train} flag, not $\tdiff$.}]
    If the mask is applied during all of training, the velocity supervision loss
    at $\tdiff = 0$ also sees a masked state---training and inference are mismatched,
    and Change~3 has no effect.
    \textit{Symptom}: adding state query input does not improve LEITMOTIF performance.
    \textit{Fix}: condition mask on \texttt{(t\_diffusion == 0)}, as in Eq.~\eqref{eq:state_mask}.

  \item[\textbf{P4: $M$ too small.}]
    High-frequency trajectory segments (contact transitions, rapid direction changes)
    are truncated.
    \textit{Symptom}: systematic position error at velocity peaks.
    \textit{Fix}: run \texttt{estimate\_M(dataset, threshold=0.99)} and increase $M$.
\end{description}

\end{document}